\documentclass[11pt]{article}
\usepackage[a4paper,margin=1in]{geometry}
\usepackage{amsmath,amssymb,bm,amsthm}
\usepackage{hyperref}

\newcommand{\rs}{r_{\mathrm{s}}}
\newcommand{\half}{\tfrac{1}{2}}

\newtheorem{proposition}{Proposition}
\newtheorem{corollary}{Corollary}

\title{Appendix 12 (to main theory): Geodesic Structure of the
  Bi-Conformal Compact Object}
\author{}
\date{}

\begin{document}
\let\phi\varphi
\maketitle

\section*{Purpose and scope}

This appendix provides the complete geodesic analysis of
the static bi-conformal metric
\begin{equation}\label{eq:metric_app12}
  ds^{2} = -e^{-\rs/r}\,c^{2}\,dt^{2}
  + e^{\rs/r}\bigl(dr^{2} + r^{2}\,d\Omega^{2}\bigr)\,,
  \qquad r > 0\,.
\end{equation}
We show that the effective potential peaks at the photon
sphere ($r = \rs$) and is exponentially suppressed to zero
as $r \to 0$.  Geodesics that cross the photon-sphere
barrier reach the curvature-regular boundary at $r = 0$ in
finite affine parameter; those that do not are reflected
and complete.  We classify this boundary and compare it
with the Schwarzschild singularity.  This supplements
Sec.~VI\,C of the main paper.


\section{Conserved quantities and geodesic equations}
\label{sec:conserved}

The metric~\eqref{eq:metric_app12} is static and spherically
symmetric.  For a geodesic $x^{\mu}(\lambda)$ with
$g_{\mu\nu}\dot{x}^{\mu}\dot{x}^{\nu} = -\kappa c^{2}$
($\kappa = 1$ timelike, $\kappa = 0$ null), the conserved
energy and angular momentum per unit mass are
\begin{align}
  \mathcal{E} &= e^{-\rs/r}\,c^{2}\,\dot{t}\,,
  \label{eq:energy}\\
  \mathcal{L} &= e^{\rs/r}\,r^{2}\,\dot{\phi}\,.
  \label{eq:ang_mom}
\end{align}
Substituting into the normalization condition and
restricting to the equatorial plane ($\theta = \pi/2$)
gives the radial equation
\begin{equation}\label{eq:radial}
  \dot{r}^{2} = \frac{\mathcal{E}^{2}}{c^{2}}
  - \kappa\,c^{2}\,e^{-\rs/r}
  - \frac{\mathcal{L}^{2}}{r^{2}}\,e^{-2\rs/r}\,.
\end{equation}
The effective potential is therefore
\begin{equation}\label{eq:Veff_app12}
  V_{\mathrm{eff}}(r) = \kappa\,c^{2}\,e^{-\rs/r}
  + \frac{\mathcal{L}^{2}}{r^{2}}\,e^{-2\rs/r}\,.
\end{equation}
Both terms vanish as $r \to 0$ due to the exponential
suppression $e^{-\rs/r}$, so
$V_{\mathrm{eff}} \to 0$ as $r \to 0^{+}$---in sharp
contrast to GR, where the centrifugal barrier
$\mathcal{L}^{2}/r^{2}$ diverges.  This is the same
effective potential used to derive the ISCO and photon
sphere in the main paper (Eq.~(31) with
$\varphi = -\rs/r$).

\paragraph{Scope of this appendix.}
The role of this appendix is to provide the geodesic
infrastructure (effective potential, photon-sphere
barrier, capture and incompleteness arguments) of the
bi-conformal exterior.  The explicit ISCO equation
$r^{2}-3\rs r+\rs^{2}=0$, its solution
$r_{\mathrm{ISCO}}=\phigr^{2}\,\rs$, and the
test-particle ISCO orbital-frequency proxy of
Sec.~``ISCO orbital-frequency proxy'' in the main paper
are
\emph{not} re-derived here; they belong to the main
text and are referenced from this appendix only.


\section{Photon-sphere barrier and capture}
\label{sec:capture}

The centrifugal term
$h(r) \equiv e^{-2\rs/r}/r^{2}$ has a unique maximum:
\begin{equation}
  \frac{dh}{dr} = \frac{2\,e^{-2\rs/r}}{r^{3}}
  \!\left(\frac{\rs}{r} - 1\right) = 0
  \qquad\Longrightarrow\qquad r = \rs\,.
\end{equation}
This is the photon sphere.  The peak value of the
centrifugal term is
$\mathcal{L}^{2}\,h(\rs) =
\mathcal{L}^{2}/(e^{2}\,\rs^{2})$.

\begin{proposition}[Finite barrier]\label{prop:barrier}
The effective potential~\eqref{eq:Veff_app12} has a finite
maximum near $r = \rs$ and vanishes as $r \to 0^{+}$.
No infinite centrifugal wall shields $r = 0$.
\end{proposition}

\begin{proposition}[Deflected geodesics]\label{prop:deflected}
Every geodesic (null or timelike) whose energy satisfies
$\mathcal{E}^{2}/c^{2} < V_{\mathrm{eff,max}}$
has a turning point at some $r_{\min} > 0$ and never
reaches $r = 0$.  Such geodesics are complete.
\end{proposition}
\begin{proof}
If $\mathcal{E}^{2}/c^{2} < V_{\mathrm{eff,max}}$,
there exists $r_{\min} > 0$ where
$\dot{r}^{2} = 0$.  For $r < r_{\min}$,
$\dot{r}^{2} < 0$, which is forbidden, so the geodesic
reverses.  It remains in $r \geq r_{\min} > 0$, where
the metric is smooth and nondegenerate.
\end{proof}

For null geodesics the critical condition is expressed
via the impact parameter $b = c\,|\mathcal{L}|/\mathcal{E}$.
A null geodesic is captured if and only if
\begin{equation}\label{eq:capture}
  b < b_{c} = \rs\,e \approx 2.718\,\rs\,,
\end{equation}
consistent with the shadow size derived in the main
paper (Sec.~VI\,E).

\begin{proposition}[Captured geodesics]\label{prop:captured}
Every geodesic (null or timelike) with
$\mathcal{E}^{2}/c^{2} > V_{\mathrm{eff,max}}$
crosses the photon-sphere barrier and reaches
$r = 0$ in finite affine parameter.
\end{proposition}
\begin{proof}
Once the geodesic passes the barrier peak, $V_{\mathrm{eff}}$
decreases monotonically to zero (since
$dV_{\mathrm{eff}}/dr > 0$ for all $r < \rs$).
Therefore $\dot{r}^{2}$ remains positive and bounded
below: $\dot{r}^{2} \geq
\mathcal{E}^{2}/c^{2} - V_{\mathrm{eff,max}} > 0$.
As $r \to 0$, $V_{\mathrm{eff}} \to 0$, so
$\dot{r}^{2} \to \mathcal{E}^{2}/c^{2}$.
The affine parameter to reach $r = 0$ is therefore bounded
by the reciprocal of the minimum radial speed,
$\Delta\lambda \leq r_{0}/|\dot{r}|_{\min} < \infty$, with
$|\dot{r}|_{\min}=\sqrt{\mathcal{E}^{2}/c^{2}-V_{\mathrm{eff,max}}}>0$.
\end{proof}

\begin{corollary}
The spacetime is geodesically incomplete.  The
incomplete geodesics comprise all trajectories that
cross the photon-sphere barrier---a positive-measure
subset of initial conditions, including all radial
geodesics ($\mathcal{L} = 0$) and all non-radial
geodesics with sufficiently small impact parameter.
\end{corollary}


\section{Radial geodesics: finite-time arrival at $r=0$}
\label{sec:radial}

For purely radial motion ($\mathcal{L} = 0$),
Eq.~\eqref{eq:radial} reduces to
\begin{equation}\label{eq:radial_simple}
  \dot{r}^{2} = \frac{\mathcal{E}^{2}}{c^{2}}
  - \kappa\,c^{2}\,e^{-\rs/r}\,.
\end{equation}

\begin{proposition}[Null radial geodesics]\label{prop:null}
Every ingoing radial null geodesic reaches $r = 0$
at finite affine parameter
$\lambda_{0} = r_{0}\,c/\mathcal{E}$.
\end{proposition}
\begin{proof}
Setting $\kappa = 0$: $\dot{r} = -\mathcal{E}/c$
(constant).  Integration gives
$r(\lambda) = r_{0} - (\mathcal{E}/c)\,\lambda$,
reaching $r = 0$ at
$\lambda_{0} = r_{0}\,c/\mathcal{E} < \infty$.
\end{proof}

\begin{proposition}[Timelike radial geodesics]\label{prop:time}
Every ingoing radial timelike geodesic
reaches $r = 0$ in finite proper time.
\end{proposition}
\begin{proof}
Setting $\kappa = 1$:
$\dot{r}^{2} = \mathcal{E}^{2}/c^{2} - c^{2}e^{-\rs/r}$.
As $r \to 0^{+}$, $e^{-\rs/r} \to 0$, so
$\dot{r}^{2} \to \mathcal{E}^{2}/c^{2} > 0$.
The proper time from $r_{0}$ to $r = 0$ is
\[
  \tau = \int_{0}^{r_{0}}
  \frac{dr}{\sqrt{\mathcal{E}^{2}/c^{2}
  - c^{2}\,e^{-\rs/r}}}\,.
\]
On the compact interval $[0,r_{0}]$ the integrand attains its
maximum at $r=r_{0}$, where $|\dot{r}|$ is minimal,
$|\dot{r}|_{\min}=\sqrt{\mathcal{E}^{2}/c^{2}-c^{2}e^{-\rs/r_{0}}}>0$
(for ingoing geodesics with $\mathcal{E}^{2}/c^{2}>c^{2}e^{-\rs/r_{0}}$,
which is automatic once the photon-sphere barrier has been crossed).
Therefore $\tau\leq r_{0}/|\dot{r}|_{\min}<\infty$.
\end{proof}


\section{Proper distance and curvature regularity}
\label{sec:regularity}

Despite the finite affine parameter, the
\emph{proper radial distance} to $r = 0$ is infinite:
\begin{equation}\label{eq:proper_distance}
  L = \int_{0}^{r_{0}} e^{\rs/(2r)}\,dr = \infty\,.
\end{equation}
The integrand diverges faster than any power of
$1/r$; the exponential factor $e^{\rs/(2r)}$
dominates all polynomial contributions.
Thus $r = 0$ is at infinite proper distance even
though geodesics reach it in finite proper
time: the proper velocity
$dL/d\tau = e^{\rs/(2r)}\,|\dot{r}| \to \infty$
diverges exponentially, while the locally measured
speed approaches $c$ from below.

For the metric~\eqref{eq:metric_app12}, the Ricci scalar, the
quadratic Ricci invariant, and the Kretschmann scalar are
\begin{equation}\label{eq:curvature}
  \begin{aligned}
  R &= -\frac{\rs^{2}}{2r^{4}}\,e^{-\rs/r}\,,\\[2pt]
  R_{\mu\nu}R^{\mu\nu}
  &= \frac{\rs^{4}}{4r^{8}}\,e^{-2\rs/r}\,,\\[2pt]
  K \equiv R_{\mu\nu\rho\sigma}R^{\mu\nu\rho\sigma}
  &= \frac{\rs^{2}}{4r^{8}}\Bigl(48r^{2}-32r\,\rs+7\rs^{2}\Bigr)\,
     e^{-2\rs/r}\,,
  \end{aligned}
\end{equation}
each tending to $0$ as $r\to 0^{+}$ faster than any inverse power of
$r$ due to the exponential factors.  Therefore $r=0$ is \emph{not a
curvature-invariant singularity}: every polynomial curvature scalar is
finite (in fact vanishes) there.  In the Hawking--Penrose sense,
however, the spacetime remains \emph{singular}, because it is
geodesically incomplete at $r=0$ (Prop.~\ref{prop:incomplete} below).
``Singularity resolution'' in this appendix refers throughout to the
curvature-invariant criterion, not to geodesic completeness.


\section{Wormhole-like topology}\label{sec:wormhole}

The areal radius is
$R(r) = r\,e^{\rs/(2r)}$.  Its derivative
\begin{equation}
  \frac{dR}{dr} = e^{\rs/(2r)}
  \!\left(1 - \frac{\rs}{2r}\right)
\end{equation}
vanishes at $r = \rs/2$, giving the throat
$R_{\min} = (\rs/2)\,e \approx 1.36\,\rs$.  For
$r > \rs/2$, $R$ increases with $r$ (exterior);
for $r < \rs/2$, $R$ increases as $r$ decreases
(interior).  Both ends open to $R \to \infty$:
\begin{equation}
  \lim_{r\to\infty} R(r) = \infty\,,
  \qquad
  \lim_{r\to 0^{+}} R(r) = \infty\,.
\end{equation}
The coordinate chart $r \in (0,\infty)$ thus covers
a geometry with two asymptotic regions
connected by a throat~\cite{Boonserm2020}.
Unlike a standard Morris--Thorne wormhole, the two
ends are physically inequivalent:
\begin{itemize}
\item \textbf{Exterior} ($r \to \infty$):
  $g_{tt} \to -c^{2}$, $g_{rr} \to 1$.
  Standard asymptotically flat region.
\item \textbf{Deep interior} ($r \to 0^{+}$):
  $g_{tt} \to 0$, $g_{rr} \to \infty$.
  Infinitely redshifted region; coordinate speed
  of light $c\,e^{-\rs/r} \to 0$ (locally measured
  speed remains $c$).
\end{itemize}


\section{Classification of the $r = 0$ boundary}
\label{sec:boundary}

\begin{proposition}\label{prop:incomplete}
The spacetime $(\mathcal{M}, g)$ with
$\mathcal{M} = \mathbb{R}\times(0,\infty)
\times S^{2}$ and metric~\eqref{eq:metric_app12} is
geodesically incomplete: all geodesics that cross the
photon-sphere barrier (including all radial geodesics)
reach $r = 0$ at finite affine parameter.
\end{proposition}

The metric function
$f(r) = e^{-\rs/r}$ extends to a $C^{\infty}$
function on $[0,\infty)$ with $f(0) = 0$ and
$f^{(n)}(0) = 0$ for all $n$ (smooth but not
real-analytic).  This observation concerns only this scalar
function, not a metric extension: in the present isotropic
coordinate chart $g_{rr}=e^{\rs/r}$ and the angular
coefficient $r^{2}e^{\rs/r}$ diverge as $r\to0$.

The bounded displayed curvature invariants therefore motivate,
but do not prove, a low-regularity metric extension.
Clarke-type extension results require additional regularity and
global hypotheses~\cite{Clarke1993,Racz2023}, and no $C^{0}$
or smooth metric extension is constructed here.  The appropriate
boundary condition or extended manifold remains open.

\medskip

\textit{Exterior perspective.}---The boundary is
unobservable from outside:
\begin{enumerate}
\item \textbf{Infinite coordinate time.}
  $\displaystyle
  t = \frac{1}{c}\int_{0}^{r_{0}}
  e^{\rs/r}\,dr = \infty$.
  No external observer registers an event at $r=0$.

\item \textbf{Infinite redshift.}
  A photon emitted at $r$ arrives at infinity with
  $1 + z = e^{\rs/(2r)} \to \infty$ as $r \to 0$.

\item \textbf{Exponential freezing.}
  Local proper time elapses as
  $d\tau/dt = e^{-\rs/(2r)} \to 0$; the deep
  interior is dynamically inert on any finite
  observation timescale.
\end{enumerate}

\textit{Infalling perspective.}---An observer in
free fall who crosses the photon-sphere barrier
reaches $r = 0$ in finite proper time
$\tau_{0} \leq r_{0}\,c/\mathcal{E}$.
Throughout the infall, the tidal acceleration is
governed by the Kretschmann scalar $K(r)$, which
rises to a finite maximum near $r\approx 0.22\,\rs$ and
then decreases monotonically toward $r=0$, vanishing
at the origin.  The observer experiences
\emph{decreasing} tidal stress as $r \to 0$---the
opposite of the Schwarzschild case, where tidal
forces diverge.  At $r = 0$ the geodesic terminates
at a curvature-free boundary; the physical
interpretation of this termination (reflecting
boundary, passage to an extended manifold, or
degenerate endpoint) remains open.

\medskip

The comparison with Schwarzschild and two other
spacetimes with regular boundaries is given in
Table~\ref{tab:comparison_app12}.

\begin{table}[ht]
\caption{Comparison of boundary properties.
\label{tab:comparison_app12}}
\begin{center}
\begin{tabular}{lcccc}
\hline\hline
 & Bi-conformal & Schwarzschild
 & Extremal RN & AdS \\
\hline
$K$ at boundary
  & $0$ & $\infty$ & finite & $0$ \\
Proper distance
  & $\infty$ & finite$^{a}$ & $\infty$ & finite \\
Proper time
  & finite & finite & finite & finite$^{b}$ \\
Observable
  & No & No & No & No \\
Curvature singularity
  & No & Yes & No & No \\
\hline\hline
\end{tabular}
\end{center}
\smallskip\noindent{\footnotesize
$^{a}$From the event horizon to $r=0$.
$^{b}$For null geodesics reaching the conformal
boundary.}
\end{table}

Present ringdown and echo searches do not yet provide a decisive
horizon test; GWTC-3 tests remain consistent with GR and find no
confirmed post-merger echoes~\cite{Cardoso2019,Miani2023,LVK2025TGR}.
The same
qualitative picture is in line with Gogberashvili's arguments
that quantum particles are reflected by the Schwarzschild
horizon and that smooth interior--exterior matching requires
material sources to lie outside the photon sphere~\cite{Gogberashvili2018,Gogberashvili2019}.


\section{Summary}

\begin{enumerate}
\item The effective potential
  $V_{\mathrm{eff}} = \kappa c^{2}e^{-\rs/r}
  + \mathcal{L}^{2}e^{-2\rs/r}/r^{2}$
  peaks at the photon sphere ($r = \rs$)
  and is exponentially suppressed to zero as
  $r \to 0$.  There is no infinite centrifugal
  barrier.

\item Geodesics that do not cross the photon-sphere
  barrier (impact parameter $b > b_{c} = \rs\,e$)
  are reflected and \textbf{complete}.  This includes
  distant flybys, stable orbits, and all
  gravitational-lensing trajectories.

\item Geodesics that cross the barrier
  ($b < b_{c}$, including all radial geodesics)
  reach $r = 0$ in finite affine parameter.
  The spacetime is \textbf{geodesically incomplete}.

\item The boundary at $r = 0$ is curvature-regular
  ($K \to 0$), at infinite proper distance, and
  at infinite redshift.  It is causally inert and
  unobservable from the exterior.

\item An infalling observer experiences
  \emph{vanishing} tidal force at $r = 0$---in sharp
  contrast to the Schwarzschild singularity, where
  tidal disruption is inevitable.

\item The capture cross-section
  ($\sigma = \pi\rs^{2}e^{2}$) is comparable to
  Schwarzschild
  ($\sigma_{\mathrm{GR}} = 27\pi\rs^{2}/4$),
  but the endpoint is qualitatively different:
  a curvature-free, unobservable boundary
  rather than a curvature singularity.
\end{enumerate}

\end{document}
